\section{Overview}
\label{sec:testing_overview}

This chapter provides a comprehensive explanation of the test structure implemented in both the backend and frontend of the Best Bike Paths system. Tests were organized following industry best practices, with clear separation between unit tests, integration tests, and end-to-end (E2E) tests.

\section{Backend Test Structure}
\label{sec:backend_test_structure}

The backend testing strategy follows a layered approach, mirroring the application architecture. Tests are organized into four main categories: \textbf{API tests}, \textbf{Infrastructure tests}, \textbf{Repository tests}, and \textbf{Service tests}.

\subsection{Test Organization}
\label{subsec:backend_test_organization}

The backend test directory is structured as follows:

The backend test suite is organized into the following modules:
\begin{itemize}
    \item \texttt{backend/tests/api/}: Tests for authentication and trip routes, ensuring correct HTTP responses and status codes.
    \item \texttt{backend/tests/infrastructure/}: Validation of server health checks and Redis connection logic.
    \item \texttt{backend/tests/repositories/}: Direct tests for data access methods in user and trip repositories.
    \item \texttt{backend/tests/services/}: Business logic validation for authentication, path management, trip recording, and weather services.
\end{itemize}

Total: 35 tests across 11 test suites are implemented to ensure backend stability.

\subsection{API Tests}
\label{subsec:api_tests}

\textbf{Location:} \texttt{backend/tests/api/}

\textbf{Purpose:} Verify that HTTP endpoints correctly handle requests, enforce authentication, and return proper status codes and data.

\subsubsection{authRoutes.test.js (3 tests)}
\label{subsubsec:auth_routes_test}

Tests the authentication endpoints using Supertest to make HTTP requests against the Express application.

\textbf{What it tests:}
\begin{itemize}
    \item \textbf{GET /api/v1/auth/me:} Validates that the authentication middleware correctly verifies JWT tokens
    \begin{itemize}
        \item Returns 401 when no token is provided
        \item Returns 403 when an invalid token is provided  
        \item Returns 200 with user data (excluding password) when a valid token is provided
        \item Returns 404 when the token is valid but user doesn't exist in database
    \end{itemize}
\end{itemize}

\textbf{Testing approach:}
\begin{itemize}
    \item Mocks database and infrastructure layers (PostgreSQL, Redis, RabbitMQ)
    \item Uses a helper function \texttt{makeToken()} to generate valid JWTs for testing
    \item Mocks \texttt{userRepository.getUserById()} to simulate database responses
    \item Verifies that sensitive data (password\_hash) is excluded from responses
\end{itemize}

\subsubsection{tripRoutes.test.js (3 tests)}
\label{subsubsec:trip_routes_test}

Tests trip management endpoints to ensure proper data retrieval and access control.

\textbf{What it tests:}
\begin{itemize}
    \item \textbf{GET /api/v1/trips/me:} Returns authenticated user's trips
    \begin{itemize}
        \item Returns 401 when no authentication token is provided
        \item Returns 200 with list of trips for authenticated user
        \item Verifies that \texttt{tripRepository.getTripsByUserId()} is called with correct user ID
    \end{itemize}
    
    \item \textbf{GET /api/v1/trips/:id:} Returns specific trip details
    \begin{itemize}
        \item Returns 200 when trip exists and belongs to authenticated user
        \item Returns 404 when trip doesn't exist
        \item Returns 404 when trip belongs to a different user (access control)
    \end{itemize}
\end{itemize}

\textbf{Testing approach:}
\begin{itemize}
    \item Mocks \texttt{tripRepository} methods to return test data
    \item Verifies authentication enforcement (401 errors)
    \item Tests ownership validation (users can only access their own trips)
    \item Uses Supertest's chaining syntax for HTTP assertions
\end{itemize}

\subsection{Infrastructure Tests}
\label{subsec:infrastructure_tests}

\textbf{Location:} \texttt{backend/tests/infrastructure/}

\textbf{Purpose:} Validate that infrastructure components (server, Redis, database connections) are properly configured and operational.

\subsubsection{server.test.js (1 test)}
\label{subsubsec:server_test}

Tests the health check endpoint to verify server availability.

\textbf{What it tests:}
\begin{itemize}
    \item \textbf{GET /health:} Returns 200 status with service information
    \begin{itemize}
        \item Response body contains \texttt{status: "ok"}
        \item Response body contains \texttt{service: "bbp-backend"}
    \end{itemize}
\end{itemize}

\textbf{Testing approach:}
\begin{itemize}
    \item Mocks Redis, RabbitMQ, and stats consumer to avoid dependencies
    \item Uses Supertest to make HTTP request to the Express app
    \item Verifies response structure and values
\end{itemize}

\subsubsection{redis.test.js (1 test)}
\label{subsubsec:redis_test}

Tests Redis connection logic.

\textbf{What it tests:}
\begin{itemize}
    \item \textbf{connectRedis():} Verifies Redis client connection
    \begin{itemize}
        \item Tests that \texttt{redisClient.isOpen} becomes true after connection
        \item Ensures connection is established successfully
    \end{itemize}
\end{itemize}

\textbf{Testing approach:}
\begin{itemize}
    \item Mocks the underlying \texttt{redis} package to avoid needing a real Redis instance
    \item Simulates connection behavior
    \item Validates wrapper logic around Redis client
\end{itemize}

\subsection{Repository Tests}
\label{subsec:repository_tests}

\textbf{Location:} \texttt{backend/tests/repositories/}

\textbf{Purpose:} Test data access layer methods that interact with the PostgreSQL database.

\subsubsection{userRepository.test.js (3 tests)}
\label{subsubsec:user_repository_test}

Tests CRUD operations for user data.

\textbf{What it tests:}
\begin{itemize}
    \item \textbf{getUserById():} Retrieves user by ID
    \begin{itemize}
        \item Returns user object when ID exists
        \item Returns undefined when ID doesn't exist
        \item Calls SQL query with correct parameters
    \end{itemize}
    
    \item \textbf{createUser():} Inserts new user
    \begin{itemize}
        \item Calls INSERT query with all user fields in correct order
        \item Handles password hash, email, username, etc.
    \end{itemize}
\end{itemize}

\textbf{Testing approach:}
\begin{itemize}
    \item Mocks the database pool's \texttt{query()} function
    \item Verifies SQL queries are called with expected parameters
    \item Uses \texttt{expect.any(String)} for flexible SQL string matching (handles whitespace variations)
    \item Validates exact parameter arrays for INSERT operations
\end{itemize}

\subsubsection{tripRepository.test.js (3 tests)}
\label{subsubsec:trip_repository_test}

Tests trip data access methods.

\textbf{What it tests:}
\begin{itemize}
    \item \textbf{getTripById():} Retrieves trip by ID
    \begin{itemize}
        \item Returns trip object when ID exists
        \item Returns undefined when ID doesn't exist
    \end{itemize}
    
    \item \textbf{createTrip():} Inserts new trip with full data
    \begin{itemize}
        \item Handles all trip fields: weather data, GPS coordinates, timestamps
        \item Returns created trip object
    \end{itemize}
\end{itemize}

\textbf{Testing approach:}
\begin{itemize}
    \item Similar structure to user repository tests
    \item Mocks database pool
    \item Validates query structure without strict string matching
\end{itemize}

\subsection{Service Tests}
\label{subsec:service_tests}

\textbf{Location:} \texttt{backend/tests/services/}

\textbf{Purpose:} Test business logic layer, validating service methods that orchestrate repository calls and implement domain rules.

\subsubsection{authService.test.js (5 tests)}
\label{subsubsec:auth_service_test}

Tests authentication logic including registration, login, and token generation.

\textbf{What it tests:}
\begin{itemize}
    \item \textbf{register():}
    \begin{itemize}
        \item Throws error if username is already taken
        \item Throws error if email is already registered
        \item Correctly hashes password using bcrypt (10 salt rounds)
        \item Creates user with email as username internally
        \item Removes password hash from returned user object
    \end{itemize}
    
    \item \textbf{login():}
    \begin{itemize}
        \item Throws "Invalid credentials" if user not found by email or username
        \item Throws "Invalid credentials" if password doesn't match
        \item Returns JWT token and user object (without password) on successful login
        \item Stores session in Redis with correct TTL
    \end{itemize}
\end{itemize}

\textbf{Testing approach:}
\begin{itemize}
    \item Mocks \texttt{userRepository}, \texttt{bcrypt}, \texttt{jsonwebtoken}, and \texttt{redisClient}
    \item Validates business rules (duplicate email/username checks)
    \item Verifies security practices (password hashing, token generation)
    \item Ensures sensitive data is excluded from responses
\end{itemize}

\subsubsection{pathService.test.js (3 tests)}
\label{subsubsec:path_service_test}

Tests path management business logic.

\textbf{What it tests:}
\begin{itemize}
    \item \textbf{addPath():} Creates new path
    \begin{itemize}
        \item Generates UUID for new paths
        \item Sets default values: userRating = 0, status = 'optimal'
        \item Calls \texttt{pathRepository.createPath()} with correct parameters
    \end{itemize}
    
    \item \textbf{updatePathStats():} Updates path rating
    \begin{itemize}
        \item Fetches average rating from \texttt{ratingRepository}
        \item Rounds rating to nearest integer
        \item Updates path with new rating
        \item Returns null if path doesn't exist
    \end{itemize}
\end{itemize}

\textbf{Testing approach:}
\begin{itemize}
    \item Mocks \texttt{pathRepository} and \texttt{ratingRepository}
    \item Verifies default value assignment
    \item Tests rating calculation logic
\end{itemize}

\subsubsection{tripService.test.js (2 tests)}
\label{subsubsec:trip_service_test}

Tests trip management logic and integration with external services.

\textbf{What it tests:}
\begin{itemize}
    \item \textbf{automaticTripStart():} Handles trip creation with weather data
    \begin{itemize}
        \item Fetches weather data from \texttt{weatherService} using GPS coordinates
        \item Creates trip with weather information (temperature, wind, rain, clouds)
        \item Rejects if weather API fails (no fallback for trip creation)
        \item Calls \texttt{tripRepository.createTrip()} with enriched trip data
    \end{itemize}
    
    \item \textbf{updateTripTrack():} Updates trip route and triggers path statistics
    \begin{itemize}
        \item Updates trip route via \texttt{tripRepository.updateTripRoute()}
        \item Only sends RabbitMQ message when \texttt{metadata.is Final = true}
        \item Does not send message for intermediate route updates
        \item Publishes to 'path\_stats\_update' queue with path ID
    \end{itemize}
\end{itemize}

\textbf{Testing approach:}
\begin{itemize}
    \item Mocks \texttt{weatherService}, \texttt{tripRepository}, \texttt{rabbitmq}, and \texttt{redisClient}
    \item Tests conditional logic (when to call external services)
    \item Verifies asynchronous message publishing behavior
\end{itemize}

\subsubsection{userService.test.js (3 tests)}
\label{subsubsec:user_service_test}

Tests user profile management logic.

\textbf{What it tests:}
\begin{itemize}
    \item \textbf{getUserProfile():}
    \begin{itemize}
        \item Retrieves user by ID from repository
        \item Removes password\_hash from returned data
        \item Returns null if user not found
        \item Throws error if userId is missing
    \end{itemize}
    
    \item \textbf{updateUserProfile():}
    \begin{itemize}
        \item Validates that userId is provided
        \item Currently throws "Not implemented" (placeholder for future feature)
    \end{itemize}
\end{itemize}

\textbf{Testing approach:}
\begin{itemize}
    \item Mocks \texttt{userRepository}
    \item Tests data sanitization (password removal)
    \item Validates input parameters
\end{itemize}

\subsubsection{weatherService.test.js (8 tests)}
\label{subsubsec:weather_service_test}

Tests external weather API integration with comprehensive coverage.

\textbf{What it tests:}
\begin{itemize}
    \item \textbf{API configuration:}
    \begin{itemize}
        \item Calls OpenWeatherMap API with correct URL
        \item Includes latitude, longitude, API key, and units='metric' in request
    \end{itemize}
    
    \item \textbf{Data mapping:}
    \begin{itemize}
        \item Correctly maps all response fields: temperature, windSpeed, windDegree, windGust, visibility, rain, clouds
        \item Handles missing \texttt{rain} field (returns 0)
        \item Handles missing \texttt{windGust} field gracefully
    \end{itemize}
    
    \item \textbf{Error handling:}
    \begin{itemize}
        \item Uses mock data when API key is missing
        \item Falls back to mock data when axios throws error
        \item Doesn't call external API if key is unavailable
    \end{itemize}
    
    \item \textbf{Singleton behavior:}
    \begin{itemize}
        \item Fixed to read \texttt{WEATHER\_API\_KEY} dynamically via \texttt{\_getApiKey()} instead of caching in constructor
        \item Allows tests to change environment variables and have them take effect immediately
    \end{itemize}
\end{itemize}

\textbf{Testing approach:}
\begin{itemize}
    \item Mocks \texttt{axios.get()} to simulate API responses
    \item Tests with various response structures (complete, partial, missing fields)
    \item Verifies fallback behavior
    \item Uses \texttt{beforeEach} to set API key dynamically
\end{itemize}

\section{Frontend Test Structure}
\label{sec:frontend_test_structure}

The frontend employs a comprehensive testing strategy covering components, screens, services, and complete user flows.

The frontend test suite is organized into a bottom-up architecture:
\begin{itemize}
    \item \texttt{frontend/\_\_tests\_\_/unit/}: Consolidated unit tests for UI components and the data layer (storage, sync queue, retry logic).
    \item \texttt{frontend/\_\_tests\_\_/integration/}: Integration tests for primary screens and feature modules (Auth, Paths, Search, Stats).
    \item \texttt{frontend/\_\_tests\_\_/flows/}: High-level end-to-end user journeys focusing on navigation and transitions.
\end{itemize}

Total: 28 tests across 5 test suites are implemented to ensure a robust frontend user experience.

\subsection{End-to-End (E2E) Tests}
\label{subsec:frontend_e2e_tests}

\textbf{Purpose:} Validate complete user journeys from UI interaction to backend API calls.

\subsubsection{authFlow.test.js}
\label{subsubsec:auth_flow_test}

Simulates the complete authentication journey.

\textbf{What it tests:}
\begin{itemize}
    \item \textbf{Registration → Login → Home flow:}
    \begin{itemize}
        \item Displays landing page with "Log in" and "Register" buttons
        \item Navigates to Login screen when "Log in" is pressed
        \item Navigates to Register screen when "Register" is pressed
        \item Submits registration form, calls API, then auto-logs in
        \item Submits login form with valid credentials, stores token in AsyncStorage
        \item Shows error message for empty email/password
        \item Shows error message for password < 6 characters
        \item Allows navigation back from Login to Home
    \end{itemize}
\end{itemize}

\textbf{Testing approach:}
\begin{itemize}
    \item Mocks backend API (\texttt{api.post}) to avoid real server dependency
    \item Mocks AsyncStorage (\texttt{setSession}, \texttt{getToken}, \texttt{getUser})
    \item Renders full navigation stack (Home, Login, Register)
    \item Uses React Navigation's \texttt{NavigationContainer} for realistic navigation testing
    \item Simulates user input with \texttt{fireEvent.changeText} and \texttt{fireEvent.press}
\end{itemize}

\subsubsection{pathFlow.test.js}
\label{subsubsec:path_flow_test}

Tests path management workflows.

\textbf{What it tests:}
\begin{itemize}
    \item \textbf{Path List → Detail → Create → Publish flow:}
    \begin{itemize}
        \item Loads and displays list of paths from \texttt{pathsApi.getPaths()}
        \item Shows "No paths" message when list is empty
        \item Navigates to Path Detail when path is pressed
        \item Displays path information: name, distance, difficulty
        \item Shows "View path data" button and navigates to data screen
        \item Allows manual path creation via PathCreateManualScreen
        \item Validates that path name is required (shows error if empty)
        \item Calls \texttt{pathsApi.createPath()} when form is submitted
        \item Displays publish screen with prefilled path information
        \item Has "Cancel" and "Back to paths" navigation options
    \end{itemize}
\end{itemize}

\textbf{Testing approach:}
\begin{itemize}
    \item Mocks \texttt{pathsApi} and \texttt{MapService}
    \item Mocks network status context
    \item Renders complete navigation stack for path screens
    \item Verifies UI state changes (loading, empty, populated)
    \item Tests form validation and submission
\end{itemize}

\subsubsection{recordTripFlow.test.js (Removed)}
\label{subsubsec:record_trip_flow_test}

Originally tested trip recording lifecycle but was removed due to complexity of sensor integration.

\textbf{Purpose (before removal):}
\begin{itemize}
    \item Starting trip recording
    \item Tracking mock GPS locations
    \item Displaying live trip progress
    \item Stopping and saving trip
\end{itemize}

\textbf{Reason for removal:} Sensor-based automatic trip detection proved too complex; manual recording is now the primary approach.

\subsubsection{syncOffline.test.js}
\label{subsubsec:sync_offline_test}

Tests offline capabilities and data synchronization.

\textbf{What it tests:}
\begin{itemize}
    \item \textbf{Offline queue management:}
    \begin{itemize}
        \item Operations queued when device is offline
        \item Queue processes automatically when connectivity is restored
        \item Error handling for failed sync attempts
        \item UI displays pending sync count
    \end{itemize}
\end{itemize}

\textbf{Testing approach:}
\begin{itemize}
    \item Simulates network status changes
    \item Mocks AsyncStorage for queue persistence
    \item Verifies retry logic
\end{itemize}

\subsection{Component Tests}
\label{subsec:component_tests}

\textbf{Purpose:} Unit tests for reusable UI components to verify rendering and event handling.

\subsubsection{UI Component Tests}
\label{subsubsec:ui_component_tests}

Tests for atomic components (Button, Card, Input, Loading, etc.):

\begin{itemize}
    \item \textbf{Button.test.js:}
    \begin{itemize}
        \item Renders with correct text
        \item Calls \texttt{onPress} handler when pressed
        \item Displays loading state when \texttt{loading=true}
        \item Applies disabled styling when \texttt{disabled=true}
    \end{itemize}
    
    \item \textbf{Card.test.js:}
    \begin{itemize}
        \item Renders children correctly
        \item Applies custom styles
        \item Handles press events for touchable cards
    \end{itemize}
    
    \item \textbf{Input.test.js:}
    \begin{itemize}
        \item Renders with placeholder text
        \item Calls \texttt{onChangeText} when text changes
        \item Displays error message when \texttt{error} prop is provided
        \item Shows/hides password with \texttt{secureTextEntry}
    \end{itemize}
    
    \item \textbf{PathCard.test.js:}
    \begin{itemize}
        \item Displays path name, distance, difficulty
        \item Renders rating stars
        \item Calls \texttt{onPress} when card is tapped
    \end{itemize}
    
    \item \textbf{PathList.test.js:}
    \begin{itemize}
        \item Renders list of PathCard components
        \item Shows loading spinner while fetching
        \item Displays empty state message when no paths
    \end{itemize}
\end{itemize}

\subsection{Screen Tests}
\label{subsec:screen_tests}

\textbf{Purpose:} Test individual application screens in isolation to ensure proper rendering and basic interactions.

\subsubsection{Primary Screen Tests}
\label{subsubsec:primary_screen_tests}

\begin{itemize}
    \item \textbf{HomeScreen.test.js:}
    \begin{itemize}
        \item Renders app title "Best Bike Paths"
        \item Shows "Log in" and "Register" buttons when unauthenticated
        \item Shows navigation tabs when authenticated
    \end{itemize}
    
    \item \textbf{LoginScreen.test.js:}
    \begin{itemize}
        \item Renders email and password input fields
        \item Renders "Enter" and "Back" buttons
        \item Shows error message for invalid input
        \item Calls login API on form submission
    \end{itemize}
    
    \item \textbf{RegisterScreen.test.js:}
    \begin{itemize}
        \item Renders registration form fields
        \item Validates email format
        \item Validates password length (minimum 6 characters)
        \item Calls register API on submission
    \end{itemize}
    
    \item \textbf{PathListScreen.test.js:}
    \begin{itemize}
        \item Fetches paths from API on mount
        \item Displays path cards
        \item Shows empty state when no paths available
    \end{itemize}
    
    \item \textbf{PathDetailScreen.test.js:}
    \begin{itemize}
        \item Displays path name, description, distance, difficulty
        \item Shows "View path data" button
        \item Shows "Publish path" button for draft paths
        \item Renders map with path route
    \end{itemize}
\end{itemize}

\subsection{Service \& Logic Tests}
\label{subsec:service_logic_tests}

\textbf{Purpose:} Test utility services and business logic on the frontend.

\subsubsection{MapService.test.js}
\label{subsubsec:map_service_test}

Tests map-related utilities:
\begin{itemize}
    \item \textbf{searchPlaces():} Mocks Mapbox geocoding API
    \item \textbf{getCurrentLocation():} Retrieves device GPS coordinates
    \item \textbf{reverseGeocode():} Converts coordinates to address
    \item \textbf{getDirections():} Fetches route between two points
    \item \textbf{formatRecordedRoute():} Converts GPS array to GeoJSON
\end{itemize}

\subsubsection{requestWithRetry.test.js}
\label{subsubsec:request_retry_test}

Tests HTTP retry logic:
\begin{itemize}
    \item Retries failed requests up to 3 times
    \item Uses exponential backoff (1s, 2s, 4s)
    \item Throws error after all retries exhausted
\end{itemize}

\subsubsection{syncQueueProcessor.test.js}
\label{subsubsec:sync_queue_processor_test}

Tests offline sync queue:
\begin{itemize}
    \item Processes queued operations in order
    \item Removes successfully completed operations
    \item Retries failed operations
    \item Handles network errors gracefully
\end{itemize}

\subsection{Storage Tests}
\label{subsec:storage_tests}

\textbf{Purpose:} Test AsyncStorage wrappers and data persistence logic.

\begin{itemize}
    \item \textbf{storageService.test.js:}
    \begin{itemize}
        \item Tests \texttt{setSession(token, user)} for saving auth data
        \item Tests \texttt{getToken()} for retrieving stored token
        \item Tests \texttt{getUser()} for retrieving user object
        \item Tests \texttt{clearSession()} for logout
    \end{itemize}
    
    \item \textbf{syncQueue.test.js:}
    \begin{itemize}
        \item Tests adding operations to queue
        \item Tests marking operations as completed
        \item Tests queue persistence across app restarts
    \end{itemize}
    
    \item \textbf{constants.test.js:}
    \begin{itemize}
        \item Validates storage key constants
        \item Ensures no key collisions
    \end{itemize}
\end{itemize}



\section{Test Execution \& Results}
\label{sec:test_execution}

\subsection{Backend Test Results}
\label{subsec:backend_test_results}

\textbf{Command:} \texttt{npm test}

\textbf{Results:}
\begin{verbatim}
Test Suites: 11 passed, 11 total
Tests:       35 passed, 35 total
Snapshots:   0 total
Time:        2.16s
\end{verbatim}

All backend tests pass successfully with 100\% pass rate.

\subsection{Frontend Test Results}
\label{subsec:frontend_test_results}

\textbf{Command:} \texttt{npm test}

\textbf{Results:}
\begin{verbatim}
Test Suites: 5 passed, 5 total
Tests:       28 passed, 28 total
Snapshots:   0 total
Time:        4.10s
\end{verbatim}

\textbf{Coverage:}
\begin{itemize}
    \item \textbf{Components \& Data:} Core UI elements and storage logic have unit tests
    \item \textbf{Screens:} Primary user-facing screens (Auth, Paths, Search, Stats) are tested
    \item \textbf{High-Level Flows:} Critical user journeys (auth, navigation) are covered
\end{itemize}

\subsection{Continuous Integration}
\label{subsec:ci_results}

Tests run automatically on:
\begin{itemize}
    \item Every commit (unit tests only)
    \item Every pull request (full test suite with Docker infrastructure)
    \item Before deployment (full suite + manual QA)
\end{itemize}

\section{Testing Best Practices Followed}
\label{sec:testing_best_practices}

\begin{itemize}
    \item \textbf{Isolation:} Each test is independent and doesn't rely on others
    \item \textbf{Mocking:} External dependencies (APIs, databases) are mocked
    \item \textbf{Descriptive Names:} Test names clearly describe what is being tested
    \item \textbf{Arrange-Act-Assert:} Tests follow AAA pattern for clarity
    \item \textbf{Test Data:} Uses realistic but non-sensitive test data
    \item \textbf{Cleanup:} Database state reset between tests (\texttt{beforeEach}/\texttt{afterEach})
    \item \textbf{Coverage:} Focus on critical paths rather than 100\% line coverage
\end{itemize}
